\infographicchapter{Endoscopy}{endoscopy.png}
\label{ch:infographic-endoscopy}

\section{Overview}

Endoscopy is a minimally invasive camera exam that lets doctors look directly inside hollow organs and body passages in real time, usually through natural openings like the mouth or rectum. Modern flexible endoscopes carry a tiny high-definition digital camera and bright light right at the tip. Paired with advanced video processing and special light filters---such as narrow-band imaging---these scopes deliver crisp, magnified, full-color views of the inner lining and its tiny blood vessels.

Doctors use endoscopy both to find problems and to treat them during the same visit across the digestive, breathing, and urinary tracts. An upper endoscopy (EGD) checks the swallowing tube (esophagus), stomach, and upper small intestine for reflux, trouble swallowing, ulcers, bleeding, or Barrett esophagus. A colonoscopy examines the large intestine to screen for colon cancer, check inflammatory bowel disease, and remove precancerous growths (polyps). Specialized types include swallowable camera capsules for the small intestine and ERCP, which combines a scope with live X-rays to diagnose and treat bile duct and pancreatic duct problems.

The biggest strengths of endoscopy are direct, true-color viewing of subtle lining changes, no X-ray radiation in standard exams (except during X-ray-assisted ERCP), and the ability to take painless tissue samples, stop bleeding, stretch narrowed areas, or remove polyps on the spot. Its limits include being an internal procedure that needs fasting or bowel prep and usually sedation, plus it only sees the inner surface unless paired with endoscopic ultrasound. Serious complications are uncommon (under 2\% for routine diagnostic exams), though small risks include sedation reactions, bleeding, a tear in the wall (perforation), infection, pancreas inflammation after ERCP, or a camera capsule getting stuck in a narrowed bowel.


\section{Underlying Physical Principles}
\begin{itemize}
\item \textbf{Bright White and Narrow-Band Light:} Powerful LED or xenon lamps shine visible white light down thin flexible light guides to brighten dark internal cavities. Special narrow-band filters select specific blue and green light colors that highlight tiny surface blood vessels and subtle tissue patterns.
\item \textbf{How Light Interacts with Tissue:} Light reflects and scatters differently depending on its color and how much blood is in the tissue. Shorter blue light stays near the very top layer to outline fine surface capillaries, while greener and redder light reaches slightly deeper vessels.
\item \textbf{Tiny Tip-Mounted Digital Camera:} Light bouncing off the organ lining passes through a miniature wide-angle lens (showing a broad 120--170 degree view) onto a tiny digital camera sensor at the tip of the scope, turning the image into crisp high-definition video.
\item \textbf{Gentle Inflation and Smart Video Tools:} Small amounts of air or quickly absorbed carbon dioxide gas gently open collapsed passages so the doctor can see clearly. Video computers enhance color and contrast, zoom in on suspicious spots, and can even run artificial intelligence (AI) tools that flag subtle polyps. Some scopes also add ultrasound (EUS) or live X-rays (ERCP) to see deeper tissues or ducts.
\end{itemize}


\section{Imaging Mechanism}
\begin{enumerate}
\item \textbf{Gentle Scope Entry and Clearing the View:} The doctor gently guides a thin, flexible video scope---which carries light guides, air/water channels, and a narrow tool channel---through a natural opening. Gentle air or carbon dioxide opens the passage while tiny water jets rinse away mucus or debris.
\item \textbf{Steering and Lighting the Lining:} Bright light illuminates the inside lining as the doctor steers the flexible tip smoothly around natural folds, switching on special narrow-band light or optical zoom whenever a closer look is needed.
\item \textbf{Capturing Digital Video:} Reflected light enters the tiny wide-angle lens at the tip of the scope and hits the digital image sensor, which converts the picture into electronic video signals sent back up the scope.
\item \textbf{Real-Time Image Enhancement:} The video processor balances colors, sharpens edges, highlights blood-vessel patterns, and---when enabled---draws live AI boxes around subtle polyps to help make sure nothing is missed.
\item \textbf{Live Screen Display and Recording:} Clear high-definition or 4K video plays instantly on room monitors to guide inspection, biopsy, or polyp removal, while key photos and clips are saved straight into your procedure report.
\end{enumerate}


\section{Key Technical Parameters}
\noindent
\begin{tabularx}{\textwidth}{@{}l X@{}}
\toprule
\textbf{Parameter} & \textbf{Typical Value / Description} \\
\midrule
Physical Principle & Direct, real-time color video viewing of internal organ linings using flexible scopes and optional narrow-band light filters \\
Energy / Signal Source & Visible white or narrow-band LED/xenon light paired with a tiny digital camera sensor at the tip of the scope \\
Typical Frequency / Energy & Visible light spectrum (400--700 nm; narrow-band mode uses blue and green light to highlight tiny surface blood vessels) \\
Spatial Resolution & High-definition detail showing fine tissue textures and tiny blood vessels down to about 10--50 micrometers with zoom \\
Typical Scan Duration & 10--20 minutes for upper endoscopy (EGD); 20--45 minutes for colonoscopy; 30--60 minutes for ERCP (plus 30--60 minutes of sedation recovery) \\
Contrast Agent & None for routine exams; occasional surface dyes to highlight tissue patterns; X-ray contrast dye used during ERCP \\
Ionizing Radiation & None for routine EGD, colonoscopy, bronchoscopy, or bladder exams; live X-rays are used during ERCP (typically 2--10 mSv) \\
\bottomrule
\end{tabularx}


\section{Clinical Usages}

\subsection{Primary Usages}
\begin{itemize}
\item \textbf{Colon Cancer Screening and Polyp Removal:} Colonoscopy is the gold-standard test for finding colon cancer early and removing precancerous growths (polyps) during the very same visit.
\item \textbf{Checking Upper Digestive Symptoms and Ulcers:} Upper endoscopy (EGD) is the main exam for trouble swallowing, stubborn acid reflux, unexplained stomach pain, ulcers, celiac disease, and monitoring Barrett esophagus.
\item \textbf{Finding and Stopping Digestive Bleeding:} Quickly locates bleeding ulcers, swollen esophageal veins (varices), or fragile surface blood vessels and treats them on the spot with tiny clips, bands, or gentle heat.
\item \textbf{Diagnosing and Tracking Inflammatory Bowel Disease:} Lets doctors see inflammation directly and take targeted tissue samples to diagnose ulcerative colitis or Crohn disease, check healing, and screen for precancerous changes.
\item \textbf{Airway and Bladder Exams:} Similar camera exams check the lungs and airways (bronchoscopy) for coughing up blood or airway blockages, and the bladder (cystoscopy) for blood in the urine or bladder growths.
\end{itemize}

\subsection{Secondary Usages}
\begin{itemize}
\item \textbf{Bile Duct and Pancreas Treatments (ERCP):} Removes trapped bile duct stones, opens blocked bile ducts with small tubes (stents), and treats leaks in the pancreatic duct.
\item \textbf{Endoscopic Ultrasound (EUS):} Pairs a camera scope with an ultrasound tip to check how deeply digestive cancers grow, look at nearby organs like the pancreas, and take targeted needle biopsies.
\item \textbf{Small-Intestine Camera Exams:} Swallowable wireless camera capsules and deep small-bowel scopes help find hidden bleeding, small-bowel Crohn disease, or inherited polyp conditions.
\item \textbf{Advanced Minimally Invasive Removal:} Specialized scope techniques (such as EMR, ESD, and POEM) can remove early digestive cancers, relieve tight swallowing muscles (achalasia), or place stents to keep narrowed passages open without open surgery.
\end{itemize}


\section{Advantages}
\begin{itemize}
\item \textbf{Direct, High-Detail Color Views:} Shows subtle redness, ulcers, fragile blood vessels, and flat precancerous spots far more clearly than outside scans like CT or X-ray.
\item \textbf{Diagnose and Treat in One Visit:} Lets doctors immediately take tissue samples, remove polyps, stop bleeding, stretch tight areas, remove bile duct stones, or place stents during the same procedure.
\item \textbf{No X-Ray Radiation in Standard Exams:} Routine upper endoscopy, colonoscopy, bronchoscopy, and cystoscopy use only visible light (only ERCP adds X-rays).
\item \textbf{Prevents Cancer and Gives Definitive Answers:} Colonoscopy actively prevents colon cancer by removing polyps before they turn malignant, and tissue samples provide exact lab answers.
\item \textbf{Gentler Alternative to Open Surgery:} Advanced scope procedures can treat early tumors and swallowing disorders with less pain, smaller risks, and much faster recovery than traditional surgery.
\end{itemize}


\section{Disadvantages}
\begin{itemize}
\item \textbf{Small Procedural Risks:} Because a scope goes inside the body, there are small risks of a tear in the organ wall (perforation), bleeding after a biopsy or polyp removal, infection, or pancreas inflammation after ERCP.
\item \textbf{Sedation and Prep Needed:} Most exams require relaxing sedation or anesthesia---which carries minor breathing or heart risks---along with fasting or a thorough bowel-cleansing prep that many people find inconvenient.
\item \textbf{Sees Only the Inner Surface:} A standard camera scope cannot see deep into the organ wall or outside organs unless paired with endoscopic ultrasound (EUS) or a scan like CT or MRI.
\item \textbf{Views Can Be Blocked:} Leftover stool from an incomplete bowel prep, tight blockages, sharp bends, or hidden creases can make small polyps harder to see; a swallowed camera capsule can also get stuck if the bowel is narrowed.
\item \textbf{Not Safe in Certain Emergencies:} Routine endoscopy is avoided if there is a suspected hole in the digestive tract, severe acute colon swelling (toxic megacolon), unstable blood pressure, or severe uncontrolled bleeding risk during high-risk treatments.
\end{itemize}


\section{Contraindications}

\subsection{Absolute}
\begin{itemize}
\item \textbf{Suspected Tears or Severe Instability:} Routine endoscopy must not be done if there is a suspected or confirmed hole in the stomach or bowel, severe life-threatening colon swelling (fulminant colitis or toxic megacolon for colonoscopy), widespread belly infection (peritonitis), unstable shock or severe breathing/heart distress, or if the patient refuses consent.
\item \textbf{Blocked Bowel (for Camera Capsules):} Swallowable wireless camera capsules must not be used if you have a known or suspected bowel narrowing, blockage, or abnormal tunnel (fistula), because the capsule could become trapped.
\end{itemize}

\subsection{Relative}
\begin{itemize}
\item \textbf{Severe Bleeding Risk:} Uncontrolled blood-clotting disorders or very low platelets make high-risk treatments like large polyp removal or duct cutting unsafe, though a simple diagnostic look may still be done if urgent.
\item \textbf{Recent Surgery, Severe Heart/Lung Disease, or Pregnancy:} Recent digestive tract surgery, a recent heart attack, severe aortic valve narrowing, or pregnancy (especially when X-rays are used in ERCP) require careful team planning, anesthesia support, or waiting if the exam is not urgent.
\end{itemize}


\section{Patient Safety Guidelines}
\begin{itemize}
\item \textbf{Safe Fasting and Bowel Prep:} Following fasting rules (typically 6--8 hours without solid food and 2 hours without clear liquids) prevents stomach contents from entering the lungs during sedation. Bowel-cleansing preps are tailored carefully for people with kidney disease, heart failure, or salt imbalances.
\item \textbf{Blood Thinners and Antibiotics:} Doctors carefully weigh bleeding risk against blood-clot risk before pausing or adjusting blood thinners. Preventive antibiotics are given only for specific higher-risk procedures, such as feeding-tube placement, bleeding varices in liver disease, or blocked bile ducts during ERCP.
\item \textbf{Close Sedation Monitoring:} Before sedation, the team checks your health and airway. Throughout the exam and recovery, staff continuously track your oxygen, heart rhythm, blood pressure, and breathing, and make sure you go home with a responsible adult.
\item \textbf{Strict Scope Cleaning:} Every scope goes through strict multi-step cleaning, leak testing, brushing, validated high-level disinfection or sterilization, forced-air drying, and tracking logs to prevent infection.
\item \textbf{Gentle Technique and Clear Warning Signs:} Gentle carbon dioxide inflation, careful safety checks for heat tools, and low-dose X-ray shielding during ERCP keep risks low. Before discharge, you receive clear written instructions on when to call for severe pain, fever, or bleeding.
\end{itemize}


\section{Prerequisite Tests}
\begin{itemize}
\item \textbf{Health and Sedation Check:} Your care team reviews your medical history, examines your airway, and checks your overall fitness for sedation, while confirming you fasted properly and completed any bowel prep.
\item \textbf{Targeted Blood Work Only When Needed:} Healthy people having a routine diagnostic scope do not need routine blood tests. Blood counts, clotting tests (PT/INR), blood-type checks, or kidney/salt panels are ordered only if you have active bleeding, liver or kidney disease, a bleeding disorder, or a higher-risk treatment planned (such as ERCP or large polyp removal).
\item \textbf{Reviewing Key Medicines:} Your doctor reviews blood thinners, aspirin or antiplatelet drugs, insulin, and GLP-1 weight-loss or diabetes medicines ahead of time to keep bleeding and sedation risks as low as possible.
\item \textbf{Pregnancy Checks and Prior Scans:} Women who could be pregnant are screened before sedation or X-ray-assisted ERCP. Before ERCP, doctors review liver labs, pancreas enzymes, and prior ultrasound, CT, or MRI scans. Before a camera capsule, a scan or dissolvable test capsule may be used if narrowing is suspected.
\end{itemize}


\section{Preparing for the Scan}
\begin{itemize}
\item \textbf{Fasting Rules:} Follow your clinic's fasting instructions closely---usually no solid food for 6--8 hours and no clear liquids for at least 2 hours before upper endoscopy, colonoscopy, or ERCP---so your stomach is empty and sedation stays safe.
\item \textbf{Colonoscopy Diet and Bowel Prep:} For a colonoscopy, you will eat low-fiber foods for a few days, switch to clear liquids (skipping red or purple drinks) the day before, and drink your prescribed split-dose cleansing solution on schedule so the doctor has a crystal-clear view.
\item \textbf{Adjusting Medications:} Ask your gastroenterologist well ahead of time about blood thinners, diabetes pills, insulin, iron pills (often paused 5--7 days before a colonoscopy), and GLP-1 medicines so they can be adjusted safely.
\item \textbf{Planning a Ride Home:} Because sedation slows your reflexes and judgment, arrange for a responsible adult to drive you home, and plan not to drive, operate machinery, or make major legal decisions for 24 hours.
\item \textbf{Clothing and Important Details:} Wear loose, comfortable clothes, leave jewelry at home, and remove dentures or glasses before an upper endoscopy. Tell the team about any implanted heart devices (important if cautery is used), drug allergies, or possible pregnancy.
\end{itemize}


\section{During Scan}
\begin{itemize}
\item \textbf{Getting Comfortable and Sedated:} You will lie comfortably on your left side (or sometimes on your back or stomach for ERCP) with oxygen, heart-rate, and blood-pressure monitors attached. IV sedation or anesthesia keeps you deeply relaxed or lightly asleep, and a small plastic mouthguard protects your teeth during an upper scope.
\item \textbf{Easy Breathing:} You breathe normally through your nose and mouth, usually with a gentle flow of extra oxygen at your nose. The flexible scope passes into the swallowing tube or colon without blocking your airway.
\item \textbf{What It Feels Like:} If you are only lightly sedated, you might notice mild throat fullness during an upper exam or brief belly bloating, cramping, or gas pressure during a colonoscopy as air or carbon dioxide gently opens the bowel. Taking biopsies or removing polyps is painless because the inner lining has no pain nerves.
\item \textbf{Constant Care:} A dedicated nurse and doctor watch your breathing, oxygen, heart rhythm, and comfort the entire time. If you are awake enough to feel sharp pain, nausea, or discomfort, raise a hand or let the team know so they can adjust right away.
\item \textbf{How Long It Takes:} Upper endoscopy usually takes 10--20 minutes, colonoscopy takes 20--45 minutes, and ERCP takes 30--60 minutes, followed by 30--60 minutes of quiet recovery as sedation wears off.
\end{itemize}


\section{Normal Appearance}
\begin{itemize}
\item \textbf{Swallowing Tube (Esophagus):} A smooth, pale-pink lining that meets the deeper orange-pink stomach lining cleanly at a neat, even border (called the Z-line), with no sores, narrowing, or swollen veins.
\item \textbf{Stomach and Upper Small Intestine:} Glistening pink stomach folds that flatten smoothly when gently inflated, a round and open exit valve (pylorus), and smooth, velvety circular folds in the upper small intestine with no ulcers or bumps.
\item \textbf{Colon (Large Intestine):} Clean, pale-pink lining showing a fine, delicate web of tiny surface blood vessels, neat natural folds, and clear landmarks at the far end of the colon, with no polyps, redness, or fragile bleeding spots.
\item \textbf{Bile and Pancreatic Ducts (During ERCP):} Smooth, gently tapering bile and pancreatic ducts of normal width (usually under 6--7 mm for the main bile duct if you still have your gallbladder), with no stones, narrowings, or dye leaks.
\end{itemize}


\section{Observing for Abnormal Patterns}
\begin{itemize}
\item \textbf{Colon Polyps and Precancerous Changes:} Using special light modes like narrow-band imaging, doctors look closely at stalked, raised, or flat growths. Darker brown color, thickened irregular tiny blood vessels, and abnormal surface patterns flag precancerous polyps or early cancer that should be removed or biopsied.
\item \textbf{Stomach Ulcers and Bleeding Risk:} Ulcers appear as distinct craters in the lining. Active spurting or oozing blood, a visible blood vessel, or a stuck clot means high risk of rebleeding and needs immediate scope treatment, while a clean white base is low risk. Thick, irregular, fragile ulcer edges raise concern for stomach cancer.
\item \textbf{Inflammatory Bowel Disease and Barrett Esophagus:} Continuous redness, lost blood-vessel patterns, easy bleeding, and shallow sores starting at the rectum point to ulcerative colitis, whereas patchy areas, deep ``cobblestone'' sores, and narrowings point to Crohn disease. Salmon-colored tongues of stomach-like tissue creeping up the lower esophagus suggest Barrett esophagus, which is biopsied and monitored.
\item \textbf{Swollen Veins and Fragile Blood Vessels:} Twisted bluish veins in the lower esophagus with red streaks are high-risk varices that can bleed heavily. Bright red, spider-like vessel clusters or red stripes in the lower stomach (``watermelon stomach'') are also common sources of slow or sudden digestive bleeding.
\item \textbf{Harmless Look-Alikes:} Gentle suction from the scope tip can leave small red marks that look like mild inflammation, and leftover mucus, bile, or stool can hide flat polyps or look like lumps until rinsed clean with water jets.
\end{itemize}


\section{Dealing with Uncertainties}
\begin{itemize}
\item \textbf{Incomplete Prep and Hidden Folds:} Leftover stool, mucus, blood, or food---as well as sharp bends and deep bowel folds---can hide flat polyps or bleeding spots and make an exam incomplete.
\item \textbf{Conditions That Look Alike:} Benign swelling bumps, irritation from acid or radiation, and poor blood flow can look very similar to precancerous changes, cancer, or Crohn disease under white light. Shallow surface biopsies can also miss deeper lumps sitting beneath the lining.
\item \textbf{Getting a Clearer Look During the Exam:} Doctors improve accuracy on the spot by rinsing thoroughly with water jets, inflating gently with carbon dioxide, switching on narrow-band light, curving the scope backward to check hidden pockets, and taking both targeted and mapping biopsies.
\item \textbf{Adding Deeper Scans:} Because standard camera scopes only see the inner lining, doctors clarify thick walls, deeper lumps, or tight narrowings with endoscopic ultrasound (EUS) and needle biopsy, CT or MRI bowel scans, or MRI duct scans (MRCP).
\item \textbf{Repeating Early or Reviewing as a Team:} If a colonoscopy prep is not clean enough for a reliable look, doctors usually recommend repeating the exam within 1 year using a stronger prep. Unclear biopsy results or bile duct narrowings are reviewed together by a team of digestive, imaging, and cancer specialists.
\end{itemize}


\section{Augmenting Methods}
\begin{itemize}
\item \textbf{Special Light Filters and Zoom:} Narrow-band and blue-light modes sharpen surface textures and tiny blood vessels in real time, making it easier to spot and judge Barrett esophagus, subtle precancerous areas, and colon polyps.
\item \textbf{Endoscopic Ultrasound (EUS) and Needle Biopsy:} Measures how deep a tumor reaches into the organ wall, checks nearby lymph nodes, examines lumps beneath the lining or in the pancreas, and guides needle biopsies for exact lab answers.
\item \textbf{CT, MRI Bowel Scans, and MRCP:} Complement camera scopes by showing deep small-bowel inflammation, tunnels (fistulas), fluid pockets outside the bowel, and noninvasive maps of the bile and pancreatic ducts.
\item \textbf{Lab Pathology and Stool Tests:} Lab examination of tissue biopsies confirms infections (like \textit{H. pylori}), celiac disease, or cancer genetics, while stool tests (FIT, stool DNA, and calprotectin) help with screening and tracking bowel inflammation.
\item \textbf{AI Polyp Detection:} Real-time artificial intelligence software watches the live video alongside the doctor, alerting them to tiny or flat colon polyps and helping tell harmless polyps apart from precancerous ones.
\end{itemize}


\section{Comparison with Alternative Modalities}
\begin{itemize}
\item \textbf{Compared with Virtual Colonoscopy (CT Colonography):} Virtual colonoscopy is noninvasive and needs no sedation, but a standard camera colonoscopy spots flat precancerous polyps more reliably and removes them right away without needing a second procedure.
\item \textbf{Compared with Barium X-Rays and CT/MRI Bowel Scans:} X-ray and CT/MRI scans check deeper bowel walls, tunnels, and outside organs without sedation, while camera endoscopy gives unmatched direct color detail and definitive tissue biopsies.
\item \textbf{Compared with MRCP (MRI of Bile Ducts):} MRCP is noninvasive, radiation-free, and preferred when doctors only need pictures of the bile and pancreatic ducts. ERCP is saved for active treatment---like removing stones or placing stents---because it carries a risk of pancreas inflammation.
\end{itemize}


\section{Follow-up Steps}
\begin{itemize}
\item \textbf{Safe Recovery and Warning Signs:} Nurses watch your vital signs until sedation wears off and release you to your adult driver. You will go home with clear instructions to seek urgent care for severe belly or chest pain, fever, vomiting blood, or black or bloody stools.
\item \textbf{Biopsy Results and Restarting Medicines:} Your clinic will share all tissue biopsy or polyp lab results with you and your regular doctor, along with clear instructions on exactly when to restart any paused blood thinners.
\item \textbf{Specialist Care if Needed:} If biopsies show cancer, advanced precancerous changes, or complex narrowings, your doctor will connect you with oncology, surgery, or advanced endoscopy specialists and arrange staging scans like CT or MRI if needed.
\item \textbf{When to Repeat the Exam:} Follow-up schedules follow proven guidelines---for instance, 10 years after a normal average-risk screening colonoscopy, 3--7 years after removing precancerous polyps (depending on size and number), 3--5 years for stable Barrett esophagus, or 8--12 weeks to confirm that a high-risk stomach ulcer has healed.
\end{itemize}


\section{What You Should Know About Endoscopy?}
\begin{itemize}
\item \textbf{Comfort and Sedation:} Endoscopy is a gentle internal camera exam done with IV sedation or anesthesia, so most people sleep comfortably and feel little to nothing---aside from a mild sore throat after an upper exam or temporary gas bloating after a colonoscopy.
\item \textbf{Radiation and Safety:} Routine upper endoscopy and colonoscopy use only visible light with zero X-ray radiation (only specialized bile-duct ERCP uses X-rays). Serious problems like bleeding or tears are uncommon, especially when fasting and bowel-prep steps are followed closely.
\item \textbf{Time and Results:} The exam itself takes 10--45 minutes plus 30--60 minutes of supervised recovery. Your doctor usually shares what they saw before you leave, while lab results from biopsies or polyps typically take 1--2 weeks.
\item \textbf{Why Your Doctor Orders It:} Endoscopy lets your doctor see ulcers, inflammation, bleeding sites, or polyps directly in sharp color, take painless tissue samples, and remove precancerous polyps to prevent colon cancer.
\item \textbf{What to Tell Your Care Team:} Follow fasting and bowel-prep instructions exactly, ask ahead about blood thinners, diabetes drugs, and GLP-1 medicines, bring a responsible adult driver, and seek urgent care afterward for severe pain, fever, or bleeding.
\end{itemize}


\section{Key References}
\begin{itemize}
\item American Society for Gastrointestinal Endoscopy (ASGE) Standards of Practice Committee. \textit{Guidelines for Sedation and Anesthesia in GI Endoscopy} and \textit{Management of Antithrombotic Agents for Patients Undergoing GI Endoscopy}. \textit{Gastrointestinal Endoscopy}. 2016--2018.
\item Gupta S, Lieberman D, Anderson JC, et al. Recommendations for follow-up after colonoscopy and polypectomy: a consensus update by the U.S. Multi-Society Task Force on Colorectal Cancer (USMSTF). \textit{Gastroenterology}. 2020;158(4):1131--1153.
\item Rex DK, Boland CR, Dominitz JA, et al. Colorectal cancer screening: recommendations for physicians and patients from the USMSTF. \textit{American Journal of Gastroenterology}. 2017;112(7):1016--1030.
\item American Society of Anesthesiologists (ASA). \textit{Practice Guidelines for Preoperative Fasting and the Use of Pharmacologic Agents to Reduce the Risk of Pulmonary Aspiration}. \textit{Anesthesiology}. 2017;126(3):376--393.
\end{itemize}
